\documentclass[11pt,a4paper]{article}
\usepackage[margin=24mm]{geometry}
\usepackage[T1]{fontenc}
\usepackage{lmodern,amsmath,amssymb}
\usepackage{microtype,fancyhdr}
\usepackage[hidelinks]{hyperref}
\hypersetup{
pdftitle={Jardim: exercise notes},
pdfsubject={Short version with sheaf primer}}
\pagestyle{fancy}
\fancyhf{}
\fancyhead[L]{\small Jardim: exercise notes}
\fancyhead[R]{\small Salvador, 2026}
\fancyfoot[C]{\thepage}
\setlength{\headheight}{14pt}
\setlength{\parindent}{0pt}
\setlength{\parskip}{4pt}
\setlength{\emergencystretch}{2em}

% latexmk -pdf -outdir=build \
%   jardim-exercises-bom.tex

\begin{document}
\thispagestyle{plain}
\begin{center}
{\Large Jardim's exercise sessions}\par
\small UFBA, Salvador,
September 28--October 1, 2026
\end{center}

\section{Sections, Rao modules and linkage}

Can two curves be joined by a chain
of links? The exercises ask for more:
can they come from sections of twists
of one common rank-two bundle?

Work on $X=\mathbb{P}^3_{\mathbb{C}}$.
Write $\mathcal{O}(a)=\mathcal{O}_X(a)$,
$F(a)=F\otimes\mathcal{O}(a)$ and
$H^i_*(F)=\bigoplus_t H^i(X,F(t))$.

For a rank-two bundle $E$ with
$\det E=\mathcal{O}(c)$, a section
$s\in H^0(E(a))$ whose zero scheme is
a curve $C$ of pure codimension two gives
\begin{equation}
0\longrightarrow\mathcal{O}(-a)
\xrightarrow{s}E
\longrightarrow\mathcal{I}_C(c+a)
\longrightarrow0.
\label{eq-section}
\end{equation}
Locally $s=(f,g)$ and $C=V(f,g)$.
Two equations of codimension two in
the smooth space $X$ form a regular
sequence: thus $C$ is locally a complete
intersection. Consequently the map
$(u,v)\mapsto fv-gu$ has image $(f,g)$
and kernel generated by $s$.
Determinants give the twist $c+a$.
The inclusion is \emph{saturated}:
its quotient is torsion-free.

\textbf{Which curves arise this way?}
Exactly the locally complete intersection
curves with $\omega_C\cong\mathcal{O}_C(m)$
for some integer $m$ (\emph{subcanonical}
curves), by Hartshorne--Serre
\cite{serre}.
Even the smooth twisted cubic fails:
$\deg\omega_C=-2$, whereas
$\deg\mathcal{O}_C(m)=3m$.

The \emph{Rao module} of a locally
Cohen--Macaulay curve is
$$
M_C=H^1_*(\mathcal{I}_C).
$$
It records sections on $C$ that do not
extend to ambient polynomials, together
with multiplication by polynomials.
The curve is arithmetically
Cohen--Macaulay precisely when $M_C=0$.

\textbf{Rao's theorem.}
For \emph{all locally Cohen--Macaulay
curves} in $X$, isomorphism of Rao modules
up to shift is equivalent to belonging
to the same \emph{even liaison class}
\cite{rao}.
A liaison is a chain of
complete-intersection links; ``even''
counts the links ($2,4,\ldots$), not
the rank of a bundle.
No presentation by a rank-two bundle
is required by the theorem.

\textbf{What a common bundle gives.}
Since $H^1(\mathcal{O}(m))=
H^2(\mathcal{O}(m))=0$ for every $m$,
two sequences \eqref{eq-section}
with the same $E$ give
$$
M_{C_1}[c+a_1]\cong H^1_*(E)
\cong M_{C_2}[c+a_2],
\qquad N[e]_t=N_{t+e}.
$$
Thus a common bundle implies even
liaison. The converse fails: a line
and a smooth conic already give a
counterexample (Section~2).

\textbf{Why introduce sections?}
They can produce the geometry of links,
not just an existence criterion.
For sections independent at the
generic point,
put $A_i=\mathcal{O}(-a_i)$ and
$T=E/(A_1+A_2)$, identifying each $A_i$
with its image. The determinant defines
a surface $S=Z(s_1\wedge s_2)$ containing
both curves. Quotienting in either order
gives
$$
\mathcal{I}_{C_1/S}(c+a_1)
\cong T\cong
\mathcal{I}_{C_2/S}(c+a_2).
$$
If $S$ is smooth, this says
$C_2\sim C_1+(a_2-a_1)H_S$, where
$H_S$ is its hyperplane class.
Choose a general $D\sim mH_S-C_1$
for $m$ large. Then both $C_1+D$ and
$C_2+D$ are hypersurface sections of
$S$: this gives $C_1\leftrightarrow
D\leftrightarrow C_2$, two links.
For $S=V(F)$ and $C_1\cup D=V(F,G)$,
the residual ideal is
$I_D=(F,G):I_{C_1}$.

\newpage
\section{A line and a smooth conic}

Let $L$ be a line and $Q$ a smooth conic.
Their ideal resolutions are
$$
\begin{aligned}
0&\to\mathcal{O}(-2)
\to\mathcal{O}(-1)^{\oplus2}
\to\mathcal{I}_L\to0,\\
0&\to\mathcal{O}(-3)
\to\mathcal{O}(-1)\oplus\mathcal{O}(-2)
\to\mathcal{I}_Q\to0.
\end{aligned}
$$
Hence both Rao modules vanish: they
are in the same even liaison class.
Individually they are zero schemes of
sections of $\mathcal{O}(1)^{\oplus2}$
and $\mathcal{O}(1)\oplus\mathcal{O}(2)$.

\textbf{A common rank-two bundle is
impossible.}
Adjunction for a regular section gives
$\omega_C\cong\mathcal{O}_C(c+2a-4)$.
On $L$ this forces $c+2a_1=2$;
on $Q$ it forces $c+2a_2=3$.
Their difference would be
$2(a_2-a_1)=1$.
The same Rao module does not enforce
the determinant needed by a common bundle.

\textbf{A common reflexive sheaf forces
coplanarity.}
Suppose instead that a rank-two
reflexive $E$ fits into
$$
\begin{aligned}
0&\to\mathcal{O}(-a_1)\xrightarrow{s_1}E
\to\mathcal{I}_L(c+a_1)\to0,\\
0&\to\mathcal{O}(-a_2)\xrightarrow{s_2}E
\to\mathcal{I}_Q(c+a_2)\to0.
\end{aligned}
$$
Put $h=c_1(\mathcal{O}(1))$.
The resolutions above give
$$
\begin{aligned}
c(\mathcal{I}_L(t))
&=1+th+h^2+(2-t)h^3,\\
c(\mathcal{I}_Q(t))
&=1+th+2h^2+(6-2t)h^3.
\end{aligned}
$$
Multiply by $1-a_i h$ and compare
the $h^2$ coefficients for $E$:
$$
1-a_1(c+a_1)=2-a_2(c+a_2),
\qquad
1=(a_2-a_1)(c+a_1+a_2).
$$
Meanwhile
$$
s_1\wedge s_2\in H^0(\mathcal{O}(d)),
\qquad d=c+a_1+a_2.
$$
This is nonzero: otherwise the two
saturated images would be the same
generic line, hence the same subsheaf,
forcing $a_1=a_2$. Thus $d\geq0$,
and the integer equation above forces
$d=1$ and $a_2-a_1=1$.
Composing either section with the other
quotient gives this same determinant
in $H^0(\mathcal{I}_L(1))$ and
$H^0(\mathcal{I}_Q(1))$, up to sign.
Therefore
$$
\boxed{L\cup Q\subset H=Z(s_1\wedge s_2).}
$$
The remaining Chern coefficients give
$c=-2a_1$, $c_2(E)=(1+a_1^2)h^2$ and
$c_3(E)=2h^3$. The last equality again
rules out local freeness. These are
necessary conditions, not an existence
proof for the common sheaf.

\textbf{How many extensions are available?}
For $C=L,Q$, Appendix~A gives
$$
\operatorname{Ext}^1
(\mathcal{I}_C(c+a),\mathcal{O}(-a))
\cong H^1(\mathcal{O}_C(c+2a-4))^\vee.
$$
The dimensions are, respectively,
$$
\max\{0,3-c-2a_1\},\qquad
\max\{0,7-2c-4a_2\}.
$$
Here $L,Q\cong\mathbb{P}^1$, but
$\mathcal{O}_Q(1)$ has degree two.
Nonzero extension classes need not give
locally free middle sheaves.

\newpage
\section{Making the maps work}

\textbf{Why is the induced map injective?}
A map $\varphi:A\to Q$ from a line
bundle to a torsion-free sheaf on an
integral variety is either zero or
injective. The local argument is short:
a map $R\to M$ sends $1$ to some $m$.
A nonzero kernel element $f$ would give
$fm=0$. Torsion-freeness forces $m=0$,
so the map is zero.

Apply this to the composite
$\varphi_1:A_1\xrightarrow{s_1}E\to Q_2$,
where $A_i=\mathcal{O}(-a_i)$ and
$Q_i=E/s_i(A_i)$ is torsion-free.
If $a_1<a_2$ and $\varphi_1=0$, then
$s_1$ factors through $A_2$, but
$$
\operatorname{Hom}(A_1,A_2)
=H^0(\mathcal{O}(a_1-a_2))=0.
$$
There are no polynomials of negative degree.
This would force $s_1=0$, a contradiction.
Thus $\varphi_1$ is injective.
For equal twists, proportional sections
must be excluded.

\textbf{Can we build the common sheaf?}
For $a_1=1$, $a_2=2$, $c=-2$, the
proposed sequences are
$$
\begin{aligned}
0&\to\mathcal{O}(-1)\to E
\to\mathcal{I}_L(-1)\to0,\\
0&\to\mathcal{O}(-2)\to E
\to\mathcal{I}_Q\to0.
\end{aligned}
$$
Both extension spaces have dimension
three. Multiplication by the plane
equation gives the cross-maps
$\mathcal{O}(-1)\to\mathcal{I}_Q$ and
$\mathcal{O}(-2)\to\mathcal{I}_L(-1)$,
with common cokernel on $H$:
$$
\mathcal{I}_{L/H}(-1)
\cong\mathcal{O}_H(-2)
\cong\mathcal{I}_{Q/H}.
$$
Recover $E$ as the fibre product
$Q_1\times_TQ_2$: pairs $(q_1,q_2)$
with the same image in $T$.
Its projections have kernels $A_1,A_2$,
giving the two sequences. Reflexivity
still requires a check.

\textbf{Where can $E$ fail to be a bundle?}
The ideal resolutions yield
$$
\begin{aligned}
\mathcal{E}xt^1(\mathcal{I}_L(t),\mathcal{O})
&\cong\mathcal{O}_L(2-t),\\
\mathcal{E}xt^1(\mathcal{I}_Q(t),\mathcal{O})
&\cong\mathcal{O}_Q(3-t).
\end{aligned}
$$
Dualizing the two sequences makes
$\mathcal{E}xt^1(E,\mathcal{O})$ a
quotient of sheaves supported on $L$
and on $Q$. Its support therefore lies
in $L\cap Q$. For reflexive $E$, that
support is exactly its non-locally-free
locus; see Appendix~A.

\textbf{Does equal cohomology suffice?}
Here is a higher-rank obstruction.
Take $F=\Omega^1_X$ and
$$
Q_1=F\oplus\mathcal{O}^{\oplus2},
\qquad
Q_2=F\oplus\mathcal{O}(1)
\oplus\mathcal{O}(-1).
$$
They have the same $H^1_*$, because
line bundles contribute nothing to it.
The Euler sequence gives $c(F)=(1-h)^4$.
Thus $c_1(Q_1)=c_1(Q_2)=-4h$, while
$c_2(Q_1)=6h^2$ and $c_2(Q_2)=5h^2$.
If both were quotients of the same $E$
with kernels $\mathcal{O}(-a_i)$,
comparison of $c_1(E)$ would force
$a_1=a_2$, and comparison of $c_2(E)$
would then be impossible.

\textbf{How might we build a bundle from $M$?}
For a finite-length graded module $M$,
choose a free presentation
$L_1\xrightarrow{\beta}L_0\to M\to0$.
Since $\widetilde M=0$, the kernel
$K=\ker\widetilde\beta$ is a bundle
with $H^1_*(K)\cong M$.
If one can additionally find a subbundle
inclusion $\alpha:\widetilde L_0^*\to K$,
then
$$
\widetilde L_0^*
\xrightarrow{\alpha}\widetilde L_1
\xrightarrow{\widetilde\beta}
\widetilde L_0
$$
is a monad with bundle cohomology
$E=K/\operatorname{im}\alpha$ and
$H^1_*(E)\cong M$.
Finding $\alpha$ is the missing step:
dualizing $\beta$ does not produce it.
Also $c_1(E)=c_1(\widetilde L_1)$,
so $c_1(E)=0$ is not automatic.

\textbf{Lefschetz question.}
Which curves have Rao modules with WLP
or SLP? For a finite-length graded
module $N$, WLP requires one linear form
$\ell$ such that $\times\ell:N_t\to N_{t+1}$
has maximal rank for every $t$.
SLP requires the same for
$\times\ell^s:N_t\to N_{t+s}$ for all
$s\geq1$. Both properties are unchanged
by grading shifts, hence are invariants
of the even liaison class. They hold
vacuously when the Rao module is zero.

\newpage
\appendix
\section{Minimal sheaf primer}

All sheaves below are coherent.
Rank and torsion are taken over an
integral variety. A sheaf is
\emph{locally free} if locally it is
$\mathcal{O}^{\oplus r}$; these are
exactly sheaves of sections of vector
bundles. Rank is dimension at the
generic point, and alone says nothing
about local freeness everywhere.

\textbf{Torsion and reflexivity.}
Over a domain, $m$ is torsion if
$fm=0$ for some nonzero scalar $f$.
Torsion-free means this forces $m=0$.
For instance $R/(x)$ has torsion over
$R=\mathbb{C}[x]$, although it has no
nonzero nilpotents.
Write $F^*=\mathcal{H}om(F,\mathcal{O}_X)$.
The canonical evaluation map is
$$
F\longrightarrow F^{**},\qquad
s\longmapsto(\varphi\mapsto\varphi(s)).
$$
It is injective precisely when $F$ is
torsion-free, and an isomorphism
precisely when $F$ is \emph{reflexive}
\cite{stacks}.
Locally free implies reflexive, which
implies torsion-free. For example,
$\mathcal{I}_L^{**}=\mathcal{O}\ne
\mathcal{I}_L$ for a line in
$\mathbb{P}^3$.

\textbf{Twisting and exactness.}
We have $F(a)(b)=F(a+b)$ and
$\operatorname{Hom}(\mathcal{O}(-a),F)
=H^0(F(a))$.
Tensoring preserves every short exact
sequence precisely for a \emph{flat}
sheaf. Locally free sheaves are flat:
if $F\cong\mathcal{O}^{\oplus r}$, then
$$
A\otimes F\cong A^{\oplus r},\qquad
a\otimes e_i\longmapsto
(0,\ldots,a,\ldots,0).
$$
Thus tensoring locally makes $r$ copies
of the sequence. In particular, twisting
by a line bundle is exact.
For arbitrary $F$, tensoring is only
right exact: the first injection may
be lost.

\textbf{Global Ext.}
$\operatorname{Ext}^1(B,A)$ classifies
sequences $0\to A\to E\to B\to0$
with end terms fixed; zero is the split
extension. Serre duality on
$\mathbb{P}^3$ gives
$$
\begin{aligned}
\operatorname{Ext}^1
(\mathcal{I}_C(c+a),\mathcal{O}(-a))
&\cong\operatorname{Ext}^2
(\mathcal{O}(-a),\mathcal{I}_C(c+a-4))^\vee\\
&\cong H^2(\mathcal{I}_C(c+2a-4))^\vee\\
&\cong H^1(\mathcal{O}_C(c+2a-4))^\vee.
\end{aligned}
$$
The second step uses the locally free
first argument; the last uses the ideal
sequence and intermediate vanishing
for $\mathcal{O}(m)$.

\textbf{Sheaf Ext: local freeness and reflexivity.}
Unlike global Ext, sheaf Ext records
local modules:
$$
\mathcal{E}xt^i(F,\mathcal{O})_x
\cong\operatorname{Ext}^i_{\mathcal{O}_{X,x}}
(F_x,\mathcal{O}_{X,x}).
$$
For torsion-free $F$ on a smooth
projective threefold, put
$\mathcal{E}^i=\mathcal{E}xt^i(F,\mathcal{O})$.
Then \cite[Section 1.1]{hl}:
\begin{itemize}
\item $F$ is locally free exactly when
$\mathcal{E}^1=\mathcal{E}^2=0$.
\item $F$ is reflexive exactly when
$\mathcal{E}^2=0$ and $\mathcal{E}^1$
has finite support.
\end{itemize}
Thus for reflexive $F$, vanishing of
$\mathcal{E}^1$ tests local freeness.
For a general torsion-free sheaf one
must also check second Ext:
$\mathcal{E}xt^1(\mathcal{I}_p,\mathcal{O})=0$
but $\mathcal{E}xt^2(\mathcal{I}_p,\mathcal{O})
\cong\mathcal{O}_p$.

\begin{thebibliography}{4}
\small\raggedright
\bibitem{serre} E. Arrondo,
\emph{A home-made Hartshorne--Serre
correspondence}, Theorem 1,
\href{https://arxiv.org/abs/math/0610015}
{arXiv:math/0610015}.
\bibitem{rao} R. Hartshorne,
\emph{On Rao's Theorems and the
Lazarsfeld--Rao Property},
\href{https://arxiv.org/abs/math/0302078}
{arXiv:math/0302078}.
\bibitem{hl} D. Huybrechts and M. Lehn,
\emph{The Geometry of Moduli Spaces
of Sheaves}, Section 1.1.
\bibitem{stacks} The Stacks Project,
\href{https://stacks.math.columbia.edu/tag/0AVT}
{Tag 0AVT}: reflexive sheaves.
\end{thebibliography}
\end{document}
